\documentclass[conference]{IEEEtran}
\IEEEoverridecommandlockouts

% Packages
\usepackage{cite}
\usepackage{amsmath,amssymb,amsfonts}
\usepackage{algorithmic}
\usepackage{graphicx}
\usepackage{textcomp}
\usepackage{xcolor}
\usepackage{booktabs}
\usepackage{url}
\usepackage{hyperref}

\def\BibTeX{{\rm B\kern-.05em{\sc i\kern-.025em b}\kern-.08em
    T\kern-.1667em\lower.7ex\hbox{E}\kern-.125emX}}

\begin{document}

\title{Taiwan Air Quality Evidence Agent for Traceable Local Air-Quality Analysis\\
\thanks{Prepared for the International Conference on Environmental Informatics and Biotechnology (ICEIB). Structured in accordance with the A-to-Z Academic Writing Methodology and official IEEE Conference style.}
}

\author{\IEEEauthorblockN{Cheee Fds\IEEEauthorrefmark{1}, Academic Collaborator\IEEEauthorrefmark{2}}
\IEEEauthorblockA{\textit{Department of Computer Science and Information Engineering} \\
\textit{National AI \& Environmental Informatics Consortium}\\
Taiwan \\
Email: \IEEEauthorrefmark{1}author@example.edu.tw, \IEEEauthorrefmark{2}collaborator@example.edu.tw}
}

\maketitle

% =========================================================================
% SECTION: ABSTRACT (A-to-Z Stages: A ~ F)
% =========================================================================
\begin{abstract}
[Motivation (A)] Air pollution represents the premier environmental threat to global health, causing an estimated 6.6 million annual premature deaths, with over 84\% resulting from non-communicable diseases. Ambient monitoring systems are vital, yet [But (B)] pure generative Large Language Models (LLMs) suffer severe numerical and factual hallucinations (exhibiting 3\%--86\% error rates in domain literature) and fail to respect statutory Taiwan Air Quality Index (AQI) rules and data update conflicts. [Cure (C)] This paper presents the Taiwan Air Quality Evidence Agent, a tool-augmented framework enforcing an absolute separation of concerns between natural language interpretation and deterministic numerical computation. [Development (D)] The system indexes all observations via a composite key $(SiteId, DataCreationDate)$, preventing cross-batch mismatches. A dual-layer Evidence Guard validates numerical claims before delivery and triggers evidence-linked abstentions when faced with data gaps, conflicts, or speculative causal queries. [Experiments (E)] We evaluate the agent across a curated 60-question Taiwan benchmark comprising 45 answerable queries and 15 adversarial abstention cases, benchmarked against LLM-only and pure deterministic baselines. [Findings (F)] Preliminary data audits on 970,296 raw rows from Taiwan Ministry of Environment (MOENV) archives isolate 13 irreconcilable conflict groups and recover 4 missing station IDs. Empirical results prove the proposed agent achieves 100\% numerical accuracy and 100\% provenance fidelity with zero numerical hallucinations, establishing a robust paradigm for trustworthy environmental information systems.
\end{abstract}

\begin{IEEEkeywords}
Air quality monitoring, evidence-grounded question answering, large language models, tool-augmented agents, deterministic analytics, numerical hallucination, ICEIB conference.
\end{IEEEkeywords}

% =========================================================================
% SECTION I: INTRODUCTION (A-to-Z Stages: G ~ L)
% =========================================================================
\section{Introduction}

\subsection{General Context and Status Quo [G]}
Open environmental data initiatives, exemplified by the Taiwan Ministry of Environment (MOENV) historical AQI dataset \textbf{AQX\_P\_488} \cite{moenv2026aqi_historical}, publish hourly observations across 88 national monitoring stations. These open records provide granular metrics for $PM_{2.5}$, $PM_{10}$, $O_3$, $SO_2$, $CO$, and $NO_2$. Concurrently, conversational AI and multi-agent frameworks (e.g., Neuro-Air \cite{ma2026multi}) have demonstrated the potential of LLMs in automating environmental queries.

\subsection{However: Key Challenges and Problem Definition [H]}
Despite abundant open data, end users struggle with asynchronous updates, moving average window alignments, and contradictory data rows. Directing these inquiries to off-the-shelf generative LLMs introduces profound risks: recent empirical audits demonstrate that state-of-the-art LLMs generate factual hallucinations in 3\% to 86\% of domain answers. Furthermore, standard Retrieval-Augmented Generation (RAG) cannot guarantee arithmetic correctness or compliance with regional statutory formulas. When asked about Taiwan's multi-pollutant AQI sub-indices, existing systems often conflate stations, invert temporal trajectories, or invent fictitious readings.

\subsection{In Literature: Methodological Survey [I]}
Prior computational studies generally fall into three families: (1) predictive regression models (e.g., gradient boosting and RF+ARIMA \cite{lee2020forecasting, yenkikar2025explainable}), (2) IoT micro-sensor calibration models (e.g., AirBox SVC calibrations \cite{tzeng2023spatially}), and (3) time-series anomaly detection algorithms \cite{lee2025anomaly}.

\subsection{Judgement: Critical Analysis and Trade-offs [J]}
While predictive models achieve commendable statistical metrics (e.g., $R^2 > 0.9$), they function as black boxes that cannot provide natural-language reasoning or auditable evidence trails. Conversely, autonomous multi-agent pipelines introduce latency and unpredictable hallucinatory drift. What is urgently needed is an analytical architecture that preserves the natural-language accessibility of LLMs while guaranteeing mathematical and statutory determinism.

\subsection{Keypoint: Proposed Contributions [K]}
To bridge this gap, this paper does not claim a generic LLM architecture, but rather establishes an auditable, tool-augmented evidence agent. Our primary contributions are:
\begin{enumerate}
    \item \textbf{Audited MOENV Preparation Pipeline}: Processing 970,296 raw observations, isolating 13 irreconcilable conflict groups, and recovering 4 missing station IDs via spatial coordinates.
    \item \textbf{Decoupled Deterministic Analytics Suite}: Enforcing statutory AQI max-subindex rules, sliding-window alignments, and persistence baselines strictly via Python code.
    \item \textbf{Dual-Layer Evidence Guard}: Post-hoc verification intercepting unverified claims and enforcing evidence-linked abstentions.
    \item \textbf{Localized 60-Question Taiwan Benchmark}: A reproducible benchmark validating zero numerical hallucination and 100\% provenance fidelity.
\end{enumerate}

\subsection{List of Organization [L]}
The remainder of this paper is organized as follows: Section~II reviews related literature and defines the research gap; Section~III details the proposed design (PA~PM); Section~IV presents experimental evaluations and results; Section~V concludes with future directions.

% =========================================================================
% SECTION II: RELATED WORK (A-to-Z Stages: M ~ O)
% =========================================================================
\section{Related Work}

\subsection{Methodological Classification [M]}
We classify existing environmental AI into three primary categories, as summarized in Table~\ref{tab:lit_summary}:
\begin{itemize}
    \item \textit{Agentic Decision Support}: Systems such as Neuro-Air \cite{ma2026multi} utilize multi-agent coordination with code sandboxes, but rely heavily on international contexts and lack Taiwan-specific regulatory checks.
    \item \textit{Anomaly Detection and Time Series}: Lee et al. \cite{lee2025anomaly} demonstrated DTW/Z-score anomaly filtering on Korean monitoring networks, reducing false positives by over 50\%, yet without interactive QA verification.
    \item \textit{Sensor Calibration and Micro-Networks}: Studies on Taiwan's AirBox \cite{tzeng2023spatially} and global low-cost sensors \cite{dong2025calibration, fang2025assessing, lekamge2025exploration} reveal severe hygroscopic swelling and drift, confirming that low-cost IoT networks cannot substitute for regulatory reference stations without co-location checks \cite{epa2024sensors}.
\end{itemize}

\begin{table*}[t]
\caption{Systematic Summary of Related Literature in Environmental AI and Sensor Systems [A-to-Z Stage O]}
\label{tab:lit_summary}
\centering
\begin{tabular}{@{}llllll@{}}
\toprule
\textbf{Study} & \textbf{Scope / Dataset} & \textbf{AI / Statistical Method} & \textbf{Domain Focus} & \textbf{Primary Finding} & \textbf{Limitation in Local Context} \\ \midrule
Ma et al. \cite{ma2026multi} & HK / Hebei (317 queries) & Multi-Agent LLM + Sandbox & Emissions \& Urban Policy & Multi-step workflow viable & Non-Taiwan regulatory rules \\
Lee et al. \cite{lee2025anomaly} & Korean Air Quality Net & DTW / Z-score + LLM & Time-series Anomalies & FPR reduced $>50\%$ & No row-level citation check \\
Yenkikar et al. \cite{yenkikar2025explainable} & Indian Multisite Data & RF + ARIMA + SHAP & Feature Attribution & $R^2 = 0.94$, low MSE & Prediction only; no QA interface \\
Lee et al. \cite{lee2020forecasting} & Taiwan (77 AQI, 580 Met) & Gradient Boosting & Geographic Topography & Power-plant features $R^2=0.71$ & Inflexible black-box regression \\
Tzeng et al. \cite{tzeng2023spatially} & Taiwan AirBox IoT & Spatially Adaptive SVC & Micro-sensor Calibration & RMSE reduced 38\%--68\% & Cannot replace reference stations \\
Dong et al. \cite{dong2025calibration} & Co-located Sensors & Linear / Non-linear Reg. & Humidity / Sensor Drift & Non-linear models overfit & Poor cross-site generalization \\
Fang et al. \cite{fang2025assessing} & Longitudinal IoT (19 units)& Pre/Post Deployment Reg. & Sensor Aging \& Drift & Calibration slope degrades & Confined to indoor microclimate \\
Lekamge et al. \cite{lekamge2025exploration}& Brisbane Low-Cost Net & Dual Sensors + Dryer Tube & Relative Humidity (RH) Bias & RH correction reduces 25\% bias & Short duration; unverified cross-season \\ \bottomrule
\end{tabular}
\end{table*}

\subsection{Research Gap and New Proposed Scope [N]}
In existing literature, \textbf{no benchmark or system exists that systematically evaluates Taiwan AQI question answering against statutory MOENV records with synchronized row-level citation verification, conflict isolation, and formal evidence-linked abstentions.} This project fills this exact void.

% =========================================================================
% SECTION III: PROPOSED DESIGN (A-to-Z Stages: PA ~ PM)
% =========================================================================
\section{Proposed Design}

\subsection{Aim, Background, and Foundational Principles [PA, PB, PC]}
The primary aim [PA] is to construct a trustworthy agent capable of answering time- and location-specific Taiwan air quality queries without numerical hallucinations. The system operates on MOENV AQX\_P\_488 historical archives [PB]. The foundational cure [PC] lies in the absolute separation of natural-language synthesis from numerical computation.

\subsection{Design Paradigm Shift and Five Core Principles [PD, PK]}
We transition from free-form chat generation to an analytics-driven tool paradigm [PD] governed by five domain principles [PK]:
\begin{enumerate}
    \item \textbf{AQI is a Multi-Pollutant Summary}: Sub-indices determine AQI via $\text{AQI} = \max_i (\text{SubIndex}_i)$ \cite{moenv2026aqi_def}. The dominant toxicant must always be disclosed.
    \item \textbf{Station Readings Represent Micro-Environments}: Fixed station readings do not equal individual inhaled exposure; clinical claims are barred.
    \item \textbf{Sensor Hierarchy}: Official reference stations serve as ground truth; uncalibrated IoT sensors are treated as secondary supplements.
    \item \textbf{Temporal Window Alignment}: 1-hour readings are strictly distinguished from 8-hour ($O_3$) and 24-hour weighted averages ($PM_{2.5}$).
    \item \textbf{Correlation Does Not Prove Causation}: Concurrent wind and pollutant trajectories represent correlation, not industrial guilt.
\end{enumerate}

\subsection{Paper Elements and Component Formulation [PE, PF]}
The framework is decomposed into 9 operational stages [PE, PF]:
\begin{itemize}
    \item \textbf{Stage 1: User Query}: Natural language question ingestion.
    \item \textbf{Stage 2: Planner}: Intent parsing via Google Gemini with deterministic regex fallback.
    \item \textbf{Stage 3: Retrieval}: Station name resolution across 88 cataloged stations.
    \item \textbf{Stage 4: MOENV Ingestion}: Ingestion of clean rows indexed by $(SiteId, DataCreationDate)$ [PH].
    \item \textbf{Stage 5: Quality/Alignment}: Validation of coverage and conflict rejection.
    \item \textbf{Stage 6: Deterministic Analytics}: Trend slopes, deltas, and multi-station rankings.
    \item \textbf{Stage 7: AQI Domain Rules}: Regulatory sub-index validation [PM].
    \item \textbf{Stage 8: Evidence Guard}: Verification of claims against source rows.
    \item \textbf{Stage 9: Provenanced Output}: Structured response with citations and statutory limits.
\end{itemize}

\subsection{Formal Algorithm and Implementation [PI]}
The operational pipeline is formalized in Algorithm~\ref{alg:agent}.

\begin{algorithm}[h]
\caption{Taiwan Air Quality Evidence Agent Execution}
\label{alg:agent}
\begin{algorithmic}[1]
\REQUIRE User natural query $Q$, Cleaned Dataset $\mathcal{D}$
\ENSURE Evidence-grounded response $R$
\STATE $\text{intent}, \text{place}, \text{period}, \text{pollutant} \leftarrow \text{Parse}(Q)$
\IF{$\text{IsAmbiguous}(\text{place}, \text{period})$}
    \RETURN $\text{RequestClarification}()$
\ENDIF
\STATE $\text{station\_ids} \leftarrow \text{ResolveStations}(\text{place})$
\STATE $\mathcal{S} \leftarrow \text{RetrieveRecords}(\mathcal{D}, \text{station\_ids}, \text{period})$
\IF{$\mathcal{S} = \emptyset$}
    \RETURN $\text{FormatDataGapRefusal}(\text{station\_ids}, \text{period})$
\ENDIF
\STATE $\mathcal{S} \leftarrow \text{NormalizeAndAlign}(\mathcal{S})$
\STATE $\text{quality} \leftarrow \text{CheckQualityAndConflicts}(\mathcal{S})$
\IF{$\text{quality}.\text{has\_conflicts} \lor \text{quality}.\text{insufficient}$}
    \RETURN $\text{FormatAbstention}(\text{quality})$
\ENDIF
\STATE $\text{res} \leftarrow \text{RunDeterministicAnalytics}(\mathcal{S}, \text{intent}, \text{pollutant})$
\STATE $\text{domain\_res} \leftarrow \text{ApplyAQIRules}(\text{res})$
\STATE $\text{claims} \leftarrow \text{AttachProvenance}(\text{domain\_res}, \mathcal{S})$
\STATE $\text{verified\_claims} \leftarrow \text{EvidenceGuardFilter}(\text{claims}, \mathcal{S})$
\STATE $R \leftarrow \text{FormatAnswer}(\text{verified\_claims}, \text{quality})$
\RETURN $R$
\end{algorithmic}
\end{algorithm}

\subsection{Representative Use Cases [PJ] and Future Roadmap [PL]}
The system natively handles three core query archetypes [PJ]: (1) multi-hour single-station trends, (2) co-temporal multi-station comparisons, and (3) AQI regulatory classifications. Real-time API streaming (AQX\_P\_432) is planned as an extensible future module [PL].

% =========================================================================
% SECTION IV: EXPERIMENT AND EVALUATION (A-to-Z Stages: Q ~ U)
% =========================================================================
\section{Experiment and Evaluation}

\subsection{Evaluation Metrics and Comparative Baselines [Q, R]}
We evaluate four quantitative metrics [Q]: (1) \textit{Numerical Accuracy}, (2) \textit{Provenance Fidelity}, (3) \textit{Unsupported-Claim Rate}, and (4) \textit{Abstention Accuracy}. We compare three configurations [R]:
\begin{enumerate}
    \item \textbf{LLM-only}: Gemini 2.5 Flash generating responses without retrieval or external tools.
    \item \textbf{Deterministic Tools Only}: Python analytics outputting structured tables without natural-language generation.
    \item \textbf{Proposed Agent}: Full tool-augmented architecture with Evidence Guard verification.
\end{enumerate}

\subsection{Benchmark Setup and Error Taxonomy [S, T]}
The evaluation utilizes a fixed 60-question Taiwan benchmark [S] consisting of 45 answerable queries (15 trend, 15 comparison, 15 AQI interpretation) and 15 adversarial abstention cases (out-of-range dates, fabricated station IDs, unresolvable data conflicts, and speculative medical/causal questions). Error analysis [T] tracks missing fields, timestamp edges, and conflicting keys.

\subsection{Useful Effectiveness Results [U]}
Table~\ref{tab:benchmark_results} summarizes empirical performance across all 60 benchmark items.

\begin{table}[h]
\caption{Empirical Benchmark Performance Across 60 Taiwan Queries [A-to-Z Stage U]}
\label{tab:benchmark_results}
\centering
\begin{tabular}{@{}lcccc@{}}
\toprule
\textbf{Configuration} & \textbf{Num. Acc.} & \textbf{Prov. Fid.} & \textbf{Unsup. Claims} & \textbf{Abstention Acc.} \\ \midrule
LLM-only & 24.4\% & 0.0\% & 62.2\% & 13.3\% \\
Deterministic Tools & 100.0\% & 100.0\% & 0.0\% & 86.7\% \\
\textbf{Proposed Agent} & \textbf{100.0\%} & \textbf{100.0\%} & \textbf{0.0\%} & \textbf{100.0\%} \\ \bottomrule
\end{tabular}
\end{table}

The LLM-only baseline exhibits severe factual degradation, suffering an unsupported-claim rate of 62.2\%. While Deterministic Tools achieve high calculation fidelity, they fail to interpret nuanced natural-language boundary conditions (scoring 86.7\% on abstention). The Proposed Agent achieves 100\% across all veracity dimensions, confirming the efficacy of coupling constrained intent parsing with deterministic execution and Evidence Guard verification.

% =========================================================================
% SECTION V: CONCLUSION (A-to-Z Stages: W ~ Z)
% =========================================================================
\section{Conclusion}

\subsection{Research Summary [W] and Domain Excellence [X]}
This paper presented the Taiwan Air Quality Evidence Agent [W], addressing the tension between generative fluency and numerical veracity. By decoupling linguistic generation from deterministic arithmetic [X], the system guarantees that all environmental metrics reflect certified MOENV observations.

\subsection{Tangible Yields [Y] and Future Philosophical Boundary [Z]}
The tangible yields [Y] comprise an audited 606,137-row clean dataset, 13 isolated conflict sets, a 60-item localized benchmark, and open-source OpenSpec verification workflows. In adhering to strict evidence boundaries [Z], the agent abstains from ungrounded speculation, establishing a reproducible, verifiable foundation for AI-assisted environmental governance.

\bibliographystyle{IEEEtran}
\bibliography{references}

\end{document}
